\section{The final obstruction and the finite cofiber}
\label{top:section}

We retain the prime $p=1009$ and put
\begin{equation}
\label{top:constants}
 N=p(p-1)-2,\qquad K=p(p-2)+1=(p-1)^2,
 \qquad W=2p^2+4p+1.
\end{equation}
The spectrum $Y$ of Proposition~\ref{low:V3} has
$BP_*Y=R/I_4$ and a binary two-sided unital multiplication
$\mu:Y\wedge Y\longrightarrow Y$ with unit $\eta:S^0_{(p)}\to Y$.
All spheres and smash products in this section are $p$-local.
We shall use the product and convergence assertions of
Proposition~\ref{fw:anss-pairing} and Proposition~\ref{fw:anss}.
In this section the second superscript of $E_r^{s,w}$ is weight:
it denotes the group of internal degree $qw$ in the earlier notation.

The May comparison and the remaining-target calculation,
Lemma~\ref{may:comparison} and Table~\ref{may:remaining-table},
give a surjection
\begin{equation}
\label{top:An-surjection}
 A^N\colon E_2^{5,W}(Y)\longrightarrow
 E_2^{qp+1,u_4+p}(Y).
\end{equation}
Indeed, every initial May monomial in the target is divisible by
$A^N$, and the two degree identities are
\[
 5+2N=qp+1,\qquad W+pN=u_4+p.
\]
For the source, $W<u_4$, so specialization to the additive Hopf
algebra identifies the normalized cobar groups in all three degrees
$4,5,6$. For the target, $u_4+p<u_5$; a positive $v_4$ coefficient
would leave weight $p$, less than the shortest adjacent tensor length
$qp=(qp+1)-1$. Thus the target groups in degrees $qp,qp+1,qp+2$
also agree with their additive specializations, by
Lemma~\ref{fw:normalized-bound}. The two triples differ by $2N$.
Cup multiplication by the actual carry cocycle representing $A^N$
commutes with specialization, so the filtered surjection of
Lemma~\ref{may:comparison} transports through this commuting square
to the displayed surjection for $Y$.

Let $\alpha\in\pi_{q-1}S^0$ and
$\beta\in\pi_{pq-2}S^0$ be the sphere classes constructed in
Lemma~\ref{top:sphere-classes}, detected by $[t_1]$ and $A$,
respectively, and write $d=|\beta|=pq-2$.
The two-cell spectrum that we use for lifts is
\begin{equation}
\label{top:J-definition}
 J=\Sigma^{-q}C(\alpha),\qquad
 \varepsilon:J\longrightarrow S^0,
\end{equation}
where $C(\alpha)$ is the cofiber of $\alpha:S^{q-1}\to S^0$
and $\varepsilon$ is its desuspended top quotient. Thus $J$ has
cells in degrees $-q$ and zero. Its connecting map is the corresponding
suspension of $\alpha$, so the cofiber exact sequence gives
\begin{equation}
\label{top:lift-criterion}
 x\in\pi_gY\text{ lifts to }\pi_g(J\wedge Y)
 \quad\Longleftrightarrow\quad
 \alpha x=0\in\pi_{g+q-1}Y.
\end{equation}

Put $n=p-1$. Since $n!$ is a $p$-local unit, the averaging idempotent
\[
 e=\frac1{n!}\sum_{\sigma\in\Sigma_n}\sigma
 \quad\text{on }J^{\wedge n}
\]
has a finite retract $T=\operatorname{Sym}^{n}J$.
Choose splitting maps $i:T\to J^{\wedge n}$ and
$\operatorname{pr}:J^{\wedge n}\to T$, with
$i\operatorname{pr}=e$, and put
$\varepsilon_T=\varepsilon^{\wedge n}i:T\to S^0$.
The key map will be
\[
 f=(\eta\beta)\circ\Sigma^d\varepsilon_T:\Sigma^dT\longrightarrow Y.
\]
The choice of $p-1$ factors has a concrete algebraic purpose:
writing $DT$ for the Spanier--Whitehead dual of $T$, we will identify
$BP_*(DT\wedge Y)$ as a free $R/I_4$-module with basis
$1,r,\ldots,r^{p-1}$ of weights $0,\ldots,p-1$.
These supply exactly the coefficient powers in the integral polynomial
$((r+x)^p-r^p-x^p)/p$.
Lemma~\ref{top:symmetric-comodule} determines the translated coaction,
and Proposition~\ref{top:symmetric-null-map} shows that substituting
$x=t_1$ gives a cochain whose cobar differential is the constant
cocycle $A$. A filtration bound then proves that $f$ is an actual
null map.

We first show that every source class is represented by
$\gamma\in\pi_{qW-5}Y$ and that
$\alpha\gamma=0$ and $\alpha\cdot(\eta\beta^p)=0$ in $\pi_*Y$.
One lift of $\gamma$ and $p-2$ lifts of $\eta\beta^p$ give, after
multiplying the $Y$ factors and averaging the $J$ factors, a map to
$T\wedge Y$ with top value $\beta^{p(p-2)}\gamma$.
The null map $f$ then gives $\beta^K\gamma=0$.
The strict inequality $5+2K<qp+1$, together with $K<N$, converts
this homotopy relation into disappearance of the whole target before
$d_{qp+1}(v_4)$.

\subsection{The small stems used for the lifts}

\begin{lemma}[Small-stem lifts]
\label{top:small-stems}
Every class of $E_2^{5,W}(Y)$ survives to homotopy. Moreover,
\begin{equation}
\label{top:alpha-gamma-vanishing}
 \pi_{q(W+1)-6}Y=0.
\end{equation}
\end{lemma}

\begin{proof}
By $q$-sparsity, a positive differential out of $(5,W)$ has target
\[
 (s,w)=(6+qj,W+j),\qquad j\geq 1.
\]
All possible Adams--Novikov terms in the stem in
\eqref{top:alpha-gamma-vanishing} have
\[
 (s,w)=(6+qj,W+1+j),\qquad j\geq 0.
\]
We treat both by writing
\begin{equation}
\label{top:small-stem-degrees}
 (s,w)=(6+qj,W+\delta+j),\qquad
 \begin{cases}
 \delta=0,&j\geq1,\\
 \delta=1,&j\geq0.
 \end{cases}
\end{equation}

First we give explicit bounds permitting the additive calculation.
The normalized-cobar inequality of Lemma~\ref{fw:normalized-bound}
implies
\[
 (q-1)j\leq W+\delta-6.
\]
At the specified prime,
\[
 q-1=2015,\quad W=2{,}040{,}199,\quad
 2015\cdot1013=2{,}041{,}195>W+1-6.
\]
Consequently $j\leq1012$ and
\begin{equation}
\label{top:small-stem-cutoff}
 w\leq2{,}041{,}212
 <p^3=1{,}027{,}243{,}729
 <u_4=1{,}028{,}262{,}820.
\end{equation}
Thus no nonconstant coefficient of $R/I_4$ occurs, in any cobar
degree at these weights. The only possible digit positions are
$0,1,2$, and the cooperation formulas are additive.
Lemma~\ref{fw:range-comparison} therefore identifies the relevant
cobar calculation with the indicated range of $P$.

Use the six exterior letters $x,a,c,y,b,z$ and the three polynomial
letters $A,B,C$ of Equation~\eqref{may:three-position-weight}.
If $m$ is the number of exterior letters and $L=(s-m)/2$ the number
of polynomial letters, then $m\leq6$ and
\[
 L\geq (p-1)j,\qquad w\geq pL.
\]
It follows that
\[
 \bigl(p(p-1)-1\bigr)j\leq W+\delta.
\]
Since $p(p-1)-1=1{,}017{,}071$ and
$3\cdot1{,}017{,}071>W+1$, this improves the bound to $j\leq2$.

For completeness, we enumerate this remaining finite calculation.
Let $H_i$ count the exterior letters covering position $i$, and let
$n_A,n_B,n_C$ be the polynomial exponents. The weight equation is
\begin{equation}
\label{top:weight-equation}
 w=H_0+p(H_1+L-n_C)+p^2(H_2+n_B+n_C).
\end{equation}
In \eqref{top:small-stem-degrees}, $m$ is even and reduction modulo
$p$ gives $H_0=1+\delta+j$, whenever a solution exists. Put
$\kappa_2=H_2+n_B+n_C$. All position contributions are nonnegative;
\eqref{top:small-stem-cutoff} also gives $w<3p^2$, so
$\kappa_2\leq2$ and $n_C\leq2$. Substituting
$L=3+(p-1)j-m/2$ into \eqref{top:weight-equation} yields
\[
 p(\kappa_2+j-2)=1+j+m/2+n_C-H_1.
\]
For $0\leq j\leq2$, its right side lies between $-3$ and $8$.
Its only multiple of $p=1009$ is zero. Hence
\begin{equation}
\label{top:small-stem-coverage}
 H_0=1+\delta+j,\qquad \kappa_2=2-j,\qquad
 H_1=1+j+m/2+n_C.
\end{equation}

For $j=2$, either $H_0=4$, which is impossible, or $H_0=3$ forces
$c$ to occur whereas $\kappa_2=0$. For $j=1,\delta=1$, all of $x,a,c$
occur and $\kappa_2=1$ excludes $b,z$. Thus $H_1\leq3$, but the last
equation of \eqref{top:small-stem-coverage}, with even $m\geq4$,
requires $H_1\geq4$. For $j=1,\delta=0$, one has $H_0=2$.
If $m=2$, the required inequality $H_1\geq3$ is impossible.
If $m=4$, one must have $H_1=4$ and $n_C=0$; the four letters
covering position one are exactly $a,c,y,b$, which have $H_2=2$,
contrary to $\kappa_2=1$. Finally $m=6$ would give $H_0=3$.
All positive-$j$ groups in \eqref{top:small-stem-degrees} therefore
have empty initial May pages.

It remains to examine $j=0,\delta=1$, where $H_0=2$ and $m$ is
$2$ or $4$. For $m=2$ both letters start at zero, and
\eqref{top:small-stem-coverage} permits only $ac$, followed by
$n_B=1,n_C=0,n_A=1$. For $m=4$ choose two of $x,a,c$ and two of
$y,b,z$. The equations become
\[
 n_C=H_1-3,\qquad n_B=2-H_2-n_C,\qquad
 n_A=1-n_B-n_C.
\]
The following table lists all nine choices; a dash means that at
least one exponent is negative.
\begin{center}
\begin{tabular}{c c c c}
\toprule
Start-zero pair & $yb$ & $yz$ & $bz$\\
\midrule
$xa$ & $xaybB$ & --- & ---\\
$xc$ & $xcybA$ & --- & ---\\
$ac$ & --- & $acyzA$ & ---\\
\bottomrule
\end{tabular}
\end{center}
Thus the complete initial-page list at $(6,W+1)$ is
\begin{equation}
\label{top:four-terms}
 acAB,\qquad xaybB,\qquad xcybA,\qquad acyzA.
\end{equation}
With exterior order $x,a,c,y,b,z$ and
$da=xy$, $db=yz$, $dc=xb+az$, their splitting differentials are
\[
\begin{aligned}
 d(acAB)&=(xab-xcy)AB, &d(xaybB)&=0,\\
 d(xcybA)&=-xaybzA, &d(acyzA)&=-xaybzA.
\end{aligned}
\]
The first image has polynomial factor $AB$ and cannot cancel either
of the other images. The two remaining cycle directions are
boundaries, since
\[
 d(acyB)=-xaybB,\qquad
 d(acbA)=-xcybA+acyzA.
\]
The splitting cohomology of \eqref{top:four-terms} is zero.
The May spectral sequence, finite in each of these weights, proves
all the claimed $E_2$ vanishings.

For orientation, the same equations also enumerate the source
$(5,W)$. Here $H_0=1$, $m$ is odd, and $L=(5-m)/2$. They force
$H_2+n_B+n_C=2$ and $H_1+L-n_C=4$. If $m=1$, then
$H_1+L\leq1+2<4$; if $m=5$, then at least two of $x,a,c$
occur, contrary to $H_0=1$. Thus $m=3$, which forces
$n_C=0$ and $H_1=3$. The only monomials are $aybB$ and $cybA$. We need no assertion that either
one survives the May spectral sequence.

No topological Adams--Novikov differential can hit filtration
five: its length would have to be at least $q+1>5$.
The outgoing target groups have just been proved zero. Thus every
actual class of $E_2^{5,W}(Y)$ survives and is represented by a
class $\gamma\in\pi_{qW-5}Y$. The vanishings for $\delta=1$ exhaust
all filtrations in the stem $q(W+1)-6$; strong convergence proves
\eqref{top:alpha-gamma-vanishing}.
\end{proof}

\subsection{Two sphere classes and an extended-power relation}

\begin{lemma}[The sphere classes]
\label{top:sphere-classes}
There are sphere classes
\[
 \alpha\in\pi_{q-1}S^0_{(p)},\qquad
 \beta\in\pi_{pq-2}S^0_{(p)}
\]
detected respectively by $[t_1]$ and the divided coproduct
$b_{1,0}$. The former generates a cyclic group of order $p$, and
\[
 \pi_iS^0_{(p)}=0\qquad(0<i<q-1).
\]
Multiplication by $\beta$ on the Adams--Novikov spectral sequence
of $Y$ induces multiplication by the actual cobar class $A$ used
in \eqref{top:An-surjection}.
\end{lemma}

\begin{proof}
The cooperation $t_1$ is primitive. In weight one the normalized
sphere cobar is
\[
 \mathbb Z_{(p)}\{v_1\}\xrightarrow{\ p\ }
 \mathbb Z_{(p)}\{t_1\},
\]
because $\eta_R(v_1)-v_1=pt_1$. Hence its degree-one cohomology is
$\mathbb Z/p$ on $[t_1]$.

The integral polynomial
\begin{equation}
\label{top:beta-cocycle}
 b(x,y)=\frac{(x+y)^p-x^p-y^p}{p}
 =\sum_{i=1}^{p-1}\frac{\binom pi}{p}x^iy^{p-i}
\end{equation}
defines the two-cobar cocycle $b_{1,0}$ by substituting the two
primitive copies of $t_1$. Its differential is zero: the
two-cocycle identity follows by expanding $(x+y+z)^p$, or by
dividing the identity $d^2(t_1^p)=0$ in the torsion-free integral
cobar. Its image in the additive specialization is nonzero in
cohomology. At weight $p$ the only coordinate available is $t_1$;
the only one-cochain of that weight is $t_1^p$, whose differential
is zero in characteristic $p$, whereas \eqref{top:beta-cocycle}
is a nonzero two-cochain. Thus it is not a boundary.

The first positive topological differential length permitted by
sparsity is $q+1$. From $(2,p)$ an outgoing differential has target
$(qj+3,p+j)$ for $j\geq1$, violating $w\geq s$; no incoming
differential can reach filtration two. The same argument applies
to $(1,1)$. Therefore these cocycles survive to the stated sphere
classes.

For any positive stem $i$ with a nonzero normalized-cobar term,
either $s=0$ and $i=qw\geq q$, or $s\geq1$ and
\[
 i=qw-s\geq(q-1)s\geq q-1.
\]
This proves the lower sphere vanishing. In stem $q-1$ the same
inequality permits only $(s,w)=(1,1)$, so there is no extension
ambiguity and $\pi_{q-1}S^0_{(p)}=\mathbb Z/p\{\alpha\}$.
Finally, the sphere action and its cobar pairing in
Proposition~\ref{fw:anss-pairing} identify the detector of
multiplication by $\beta$ with $b_{1,0}=A$.
\end{proof}

We next prove that $\alpha\beta^p$ has zero image in $\pi_*Y$. We use
the odd-primary reduced powers, including the Cartan formula and
$\mathcal P^k(x)=x^p$ for a space class of degree $2k$, in
\citet[Chapter~VI, Section~1, and Chapter~VIII, Theorem~2.2]{SteenrodEpstein1962}.
Their compatibility with suspension gives operations on spectra; the
instability normalization below is applied to a Thom space before
desuspension. Toda proved the corresponding sphere relation
\citep[p.~839, Corollary]{Toda1967} and subsequently treated such
relations using cyclic extended powers
\citep[Sections~1--3, especially Theorem~3]{Toda1968}.
We give the image relation needed in $Y$ through the following
finite-cell construction.

\begin{lemma}[A finite three-cell model]
\label{top:three-cell-model}
Let $U_0$ be a connective $p$-local spectrum with
\[
 H_0(U_0;\mathbb Z_{(p)})=\mathbb Z_{(p)},\qquad
 H_{q-1}(U_0;\mathbb Z_{(p)})=\mathbb Z/p,
\]
and with all other positive integral homology through degree $q$
zero. Let $b_0:S^0\to U_0$ induce an isomorphism on $H_0$, and
suppose the reduced power $\mathcal P^1$ of the bottom mod-$p$
cohomology class is nonzero. Then there is a three-cell spectrum
\[
 Q=\bigl(S^0\vee S^{q-1}\bigr)
       \cup_{(a\alpha,\,up)}e^q
\]
and a map $Q\to U_0$ whose restriction to $S^0$ is $b_0$.
Here $u\in\mathbb Z_{(p)}^\times$ and
$a\in\mathbb F_p^\times$.
\end{lemma}

\begin{proof}
Let $C$ be the cofiber of $b_0:S^0\to U_0$.
The assumed integral homology groups and successive applications
of connective Hurewicz show that $C$ is $(q-2)$-connected and
\[
 \pi_{q-1}C=H_{q-1}(C;\mathbb Z_{(p)})=\mathbb Z/p.
\]
Choose a generator $S^{q-1}\to C$. Its obstruction to lifting to
$U_0$ lies in $\pi_{q-2}S^0$, which is zero by
Lemma~\ref{top:sphere-classes}. A lift, together with $b_0$, gives
$S^0\vee S^{q-1}\to U_0$.
Its cofiber $D$ is $(q-1)$-connected. The integral homology exact
sequence gives
\[
 H_q(D;\mathbb Z_{(p)})
 =\ker\bigl(\mathbb Z_{(p)}\longrightarrow\mathbb Z/p\bigr)
 =p\mathbb Z_{(p)}.
\]
In particular this is a free rank-one module.
Use Hurewicz to choose a generator of $\pi_qD$ and take its
boundary in $\pi_{q-1}(S^0\vee S^{q-1})$. Attaching the
corresponding cell gives a map $Q\to U_0$ extending the two lower
cells. Its component on $S^{q-1}$ is multiplication by a unit
times $p$, since its homology image is exactly
$p\mathbb Z_{(p)}$. Its component on the bottom belongs to
$\pi_{q-1}S^0=\mathbb Z/p\{\alpha\}$, so it is $a\alpha$ for
some $a\in\mathbb F_p$.

The map $Q\to U_0$ induces isomorphisms on integral homology in
degrees zero and $q-1$. The universal coefficient sequence then
gives isomorphisms on mod-$p$ cohomology in degrees zero and $q$;
the latter is the class arising from the torsion in degree $q-1$.
The assumed nonzero reduced power therefore pulls back to a
nonzero reduced power of the bottom class of $Q$.
If $a=0$, the actual attaching map would give a wedge
\[
 Q\simeq S^0\vee\operatorname{cofib}
       \bigl(up:S^{q-1}\to S^{q-1}\bigr).
\]
The degree-zero class in this wedge is pulled back from $S^0$,
so its positive reduced powers vanish by naturality. This is a
contradiction. Hence $a\ne0$.
The construction used an integral homology gap and relative
Hurewicz in a finite range; no assertion that $U_0$ itself is
finite is involved.
\end{proof}

\begin{proposition}[The extended-power relation]
\label{top:extended-power-relation}
The image of $\alpha\beta^p$ in $\pi_*Y$ is zero.
\end{proposition}

\begin{proof}
Write $d=pq-2=2m$, where
\[
 m=p(p-1)-1\equiv-1\pmod p.
\]
Let $\rho_{\mathbb C}$ be the complex regular representation of
$C_p$. The cyclic extended power of the even sphere $S^{2m}$ is
the Thom spectrum of $m\rho_{\mathbb C}$ over $BC_p$:
\[
 (EC_p)_+\wedge_{C_p}(S^{2m})^{\wedge p}
 \simeq\operatorname{Th}(m\rho_{\mathbb C}).
\]
Applying this functor to a sphere map representing $\beta$, then
collapsing $BC_p$ in the extended power of $S^0$, gives a map
\begin{equation}
\label{top:extended-beta-map}
 \Phi:\operatorname{Th}(m\rho_{\mathbb C})\longrightarrow S^0
\end{equation}
whose restriction to a fiber sphere $S^{pd}$ is $\beta^p$.
Only the symmetric multiplication of the sphere spectrum is used
in constructing this map.

The group $\mathbb F_p^\times$ acts by permutations of the
cyclically labelled smash coordinates and by the corresponding
automorphisms of $C_p$. This normalizer action is compatible with
\eqref{top:extended-beta-map}, on whose target it is trivial.
Normalize the bottom degree by setting
\[
 U=\Sigma^{-pd}\operatorname{Th}(m\rho_{\mathbb C}),\qquad
 e_0=\frac1{p-1}\sum_{a\in\mathbb F_p^\times}a:U\longrightarrow U.
\]
The idempotent $e_0$ splits in spectra. Denote its image by $U_0$,
with maps $i_0:U_0\to U$, $r_0:U\to U_0$ satisfying
$i_0r_0=e_0$ and $r_0i_0=1$.
The projected Thom fiber $b_0:S^0\to U_0$ induces an isomorphism
on $H_0(-;\mathbb Z_{(p)})$: the complex orientation makes every
normalizer permutation act as one on the bottom class.
Moreover, normalizer invariance of $\Phi$ gives a map
\[
 \Phi_0:\Sigma^{pd}U_0\longrightarrow S^0,
 \qquad \Phi_0\Sigma^{pd}b_0=\beta^p.
\]

We describe the homology and the reduced power of this summand.
The complex bundle is integrally oriented, so the integral Thom
isomorphism identifies $H_*(U;\mathbb Z_{(p)})$ with
$H_*(BC_p;\mathbb Z_{(p)})$. The periodic cyclic-group resolution,
with maps $g-1$ and $1+g+\cdots+g^{p-1}$, gives
\[
 H_0(BC_p;\mathbb Z_{(p)})=\mathbb Z_{(p)},\quad
 H_{2j-1}(BC_p;\mathbb Z_{(p)})=\mathbb Z/p\ (j\geq1),\quad
 H_{2j}(BC_p;\mathbb Z_{(p)})=0\ (j\geq1).
\]
These groups are finitely generated degreewise, as are their
direct summands for $U_0$. The same resolution, or the circle
bundle of Chern class $p$ over infinite complex projective space,
gives
\[
 H^*(BC_p;\mathbb F_p)=\Lambda(s)\otimes\mathbb F_p[t],
 \qquad |s|=1,\quad |t|=2,\quad t=\operatorname{Bockstein}(s).
\]
An automorphism $a\in\mathbb F_p^\times$ multiplies both $s$ and
$t$ by $a$. Equivalently it acts by $a^j$ on integral homology
in degree $2j-1$. Averaging consequently retains, through degree
$q$, exactly the mod-$p$ classes in degrees $0,q-1,q$ and the
integral groups
\begin{equation}
\label{top:U0-homology}
 H_0(U_0;\mathbb Z_{(p)})=\mathbb Z_{(p)},\quad
 H_{q-1}(U_0;\mathbb Z_{(p)})=\mathbb Z/p,
 \quad H_j(U_0;\mathbb Z_{(p)})=0\ (0<j\leq q,\ j\ne q-1).
\end{equation}

Let $u$ be the normalized Thom class. We claim
\begin{equation}
\label{top:Thom-P1}
 \mathcal P^1u=-m\,t^{p-1}u\ne0.
\end{equation}
For a universal complex line bundle with first Chern class $l$
and Thom class $u_l$, write
$\mathcal P^1u_l=c\,l^{p-1}u_l$ by the Thom isomorphism.
Pullback along the zero section and the identity
$\mathcal P^1(l)=l^p$ imply $c=1$ in the polynomial cohomology of
infinite complex projective space. Naturality gives the formula
for every line bundle. The Cartan formula for a direct sum of
lines therefore gives the sum of their $(p-1)$st Chern powers as
the multiplier of $\mathcal P^1$ on its Thom class.
The regular representation splits into the characters with first
Chern classes $it$, $i\in\mathbb F_p$. Thus the multiplier here is
\[
 m\sum_{i\in\mathbb F_p}(it)^{p-1}
 =-m\,t^{p-1}.
\]
Both this class and $u$ are normalizer-invariant. The reduced
power commutes with the idempotent, so the same nonzero formula
holds in $U_0$.

Apply Lemma~\ref{top:three-cell-model} to
\eqref{top:U0-homology} and \eqref{top:Thom-P1}. It supplies a
three-cell spectrum
\[
 Q=(S^0\vee S^{q-1})\cup_{(a\alpha,\,up)}e^q\longrightarrow U_0
\]
with $a$ nonzero modulo $p$ and $u$ a $p$-local unit.

Restrict $\Phi_0$ to $\Sigma^{pd}Q$. Its bottom value is
$\beta^p$; let its value on the other lower cell be $\zeta$.
The top cell gives the relation
\[
 a\alpha\beta^p+up\zeta=0
 \quad\text{in }\pi_{pd+q-1}S^0.
\]
The action of $p$ on $Y$ is null by Corollary~\ref{low:p-null}.
Apply the unit to this relation and use $a\ne0$ modulo $p$ to
obtain $\alpha\beta^p=0$ in $\pi_*Y$.
\end{proof}

\subsection{The symmetric comodule and the null map}

We now prove that the map $f:\Sigma^dT\to Y$ defined above is null.
Finite duality turns it into the bottom image of $\beta$ in
$\pi_d(DT\wedge Y)$. The next lemma identifies the coefficient
comodule in which we will make this image a cobar boundary.

\begin{lemma}[The symmetric coefficient comodule]
\label{top:symmetric-comodule}
For $T=\operatorname{Sym}^{p-1}(\Sigma^{-q}C(\alpha))$, the module
$BP_*(DT\wedge Y)$ is free over $R/I_4$ with basis
$1,r,\ldots,r^{p-1}$, where $r^k$ has weight $k$. Its coaction is
\begin{equation}
\label{top:translated-comodule}
 r^k\longmapsto\sum_{j=0}^k\binom kjr^jt_1^{k-j}.
\end{equation}
The induced homology map of
$D\varepsilon_T\wedge1_Y:Y\to DT\wedge Y$ is the inclusion
of constants. The powers of $r$ denote basis elements; no
multiplication on $DT$ is asserted.
\end{lemma}

\begin{proof}
\smallskip\noindent\emph{The two-cell coaction.}
We first compute the coefficient comodule needed for duality.
The dual $DJ$ is the cofiber of
$\alpha':S^{q-1}\to S^0$, where $\alpha'=\pm\alpha$ after fixing
cell orientations, and its bottom inclusion is $D\varepsilon$.
Put $\mathcal F=\operatorname{fib}(S^0\xrightarrow{\eta_{BP}}BP)$.
A chosen $BP$-nullhomotopy of $\alpha'$ gives a lift
$\widetilde\alpha:S^{q-1}\to\mathcal F$ and a map $H:DJ\to BP$
extending the unit, in the following map of cofiber sequences:
\[
\begin{array}{ccccccc}
 S^{q-1}&\xrightarrow{\alpha'}&S^0&\longrightarrow&DJ&
       \longrightarrow&S^q\\[2pt]
 \big\downarrow\mathrlap{\,\scriptstyle\widetilde\alpha}&&
 \big\Vert&&\big\downarrow\mathrlap{\,\scriptstyle H}&&
 \big\downarrow\mathrlap{\,\scriptstyle\Sigma\widetilde\alpha}\\[2pt]
 \mathcal F&\longrightarrow&S^0&\xrightarrow{\eta_{BP}}&BP&
       \longrightarrow&\Sigma\mathcal F .
\end{array}
\]
The $BP$-module $BP\wedge DJ$ splits because $BP_{q-1}=0$.
Let $e_0$ be its bottom generator and choose an upper generator
$z\in BP_qDJ$ whose top quotient is $1$ and for which
$\rho_*z=0$, where $\rho=\mu_{BP}(BP\wedge H):BP\wedge DJ\to BP$
retracts the bottom inclusion. With $\epsilon:\Gamma\to R$
the counit, the element $g=(BP\wedge H)_*z$ satisfies
\[
 \epsilon(g)=0,\qquad
 \Gamma_q=\eta_L(R_q)\oplus\mathbb Z_{(p)}\{t_1\},
 \qquad g=ct_1.
\]
Here $\eta_L$ is induced by the bottom map
$BP\wedge\eta_{BP}:BP\to BP\wedge BP$.
After smashing the cofiber diagram with $BP$, its connecting
map $\partial:\Gamma_q\to\pi_{q-1}(BP\wedge\mathcal F)$ has kernel
$\eta_L(R_q)$ and gives
\[
 \partial g=(BP\wedge\widetilde\alpha)_*(1).
\]
Thus $g$ represents the first Adams--Novikov detector of
$\alpha'$. It is also the actual off-diagonal coaction:
writing $\psi(z)=1\otimes z+\omega\otimes e_0$, naturality
under $H$ gives
$\Delta g=1\otimes g+\omega\otimes1$.
Apply $\operatorname{id}\otimes\epsilon$ and use
$\epsilon(g)=0$ to obtain $\omega=g$.
By Lemma~\ref{top:sphere-classes} this detector is
$\pm[t_1]$, so $c\not\equiv0\pmod p$.

For completeness, fix the left-comodule convention
$\psi(am)=\eta_L(a)\psi(m)$ with tensor product using $\eta_R$,
so $\eta_R(v_1)-\eta_L(v_1)=pt_1$.
Replacing $z$ by $z+av_1e_0$ changes its off-diagonal term to
$g+a(\eta_L-\eta_R)(v_1)=g-pa\,t_1$.
Changing $H$ by the composite
$DJ\to S^q\xrightarrow{av_1}BP$ instead changes $g$ by
$\eta_R(av_1)$; renormalizing the upper generator to
$z-av_1e_0$ subtracts $\eta_L(av_1)$, giving $g+pa\,t_1$.
Opposite conventions reverse these signs, and dual cell
orientations change the detector by a sign; none changes its
nonzero residue modulo $p$. The underlying $BP$-module splitting
therefore need not split the comodule.
Replacing $z$ by $c^{-1}z$ now makes the off-diagonal term
exactly $t_1$. Its top quotient is then $c^{-1}$ rather than $1$;
we no longer need the earlier top normalization. This change fixes
$e_0$ and the bottom inclusion $D\varepsilon$, which are the data
used in the following symmetric-power calculation.

\smallskip\noindent\emph{The symmetric coefficient module.}
Duality changes each permutation into its inverse, and hence
preserves the average. Consequently
\[
 DT\simeq\operatorname{Sym}^{n}(DJ).
\]
Under this identification $D\varepsilon_T$ is the pure bottom
tensor $e_0^{\otimes n}$. All the cell degrees of $DJ$ are even.
If $z_S$ is the tensor with upper generator in positions $S$ and
bottom generator elsewhere, define its normalized symmetric
basis by
\begin{equation}
\label{top:symmetric-basis}
 r_k=\binom nk^{-1}\sum_{|S|=k}z_S,
 \qquad 0\leq k\leq n.
\end{equation}
Every denominator is a $p$-local unit. Expanding the translation
of each tensor gives the coefficient
\[
 \frac{\binom{n-j}{k-j}\binom nj}{\binom nk}=\binom kj
\]
on $r_jt_1^{k-j}$ in the coaction of $r_k$.
Thus we can denote this basis by $1,r,\ldots,r^{p-1}$, with
coaction~\eqref{top:translated-comodule}. Each $r^k$ has weight $k$
because it is represented by tensors with $k$ upper generators of
degree $q$.

The module $BP_*DT$ is free. Smashing its free $BP$-module
splitting with $Y$ gives
\[
 BP_*(DT\wedge Y)=BP_*DT\otimes_R(R/I_4),
\]
with the same translated basis and with $D\varepsilon_T$ identified
with constants. This argument uses a free factor, so it requires
no flatness of $R/I_4$.
\end{proof}

\begin{proposition}[The symmetric null map]
\label{top:symmetric-null-map}
For $d=|\beta|=pq-2$, the actual stable map
\begin{equation}
\label{top:null-composite}
 f:\Sigma^dT\xrightarrow{\ \Sigma^d\varepsilon_T\ }S^d
       \xrightarrow{\ \beta\ }S^0
       \xrightarrow{\ \eta\ }Y
\end{equation}
is nullhomotopic.
\end{proposition}

\begin{proof}
Finite duality identifies
\[
 [\Sigma^dT,Y]\cong[S^d,DT\wedge Y]
\]
and carries $f$ to
$(D\varepsilon_T\wedge1_Y)(\eta\beta)$.
By Lemma~\ref{top:symmetric-comodule}, its induced $E_2$ class is
represented by the constant-valued beta cocycle. We first show that this cocycle
is a boundary in the translated coefficient module.
Consider the integral polynomial
\begin{equation}
\label{top:carry-polynomial}
 F(r,x)=\frac{(r+x)^p-r^p-x^p}{p}
       =\sum_{i=1}^{p-1}\frac{\binom pi}{p}r^{p-i}x^i.
\end{equation}
The displayed sum defines it before reduction modulo $p$; it uses
only the allowed coefficient powers $r,\ldots,r^{p-1}$.
The exact integral identity
\begin{equation}
\label{top:carry-boundary}
 F(r,x)+F(r+x,y)-F(r,x+y)
       =\frac{(x+y)^p-x^p-y^p}{p}
\end{equation}
follows by cancellation of the four $p$th powers.
After reduction, the coefficient coaction
\eqref{top:translated-comodule} and primitivity of $t_1$ make the
left side the cobar differential of $F(r,t_1)$, with the common
overall sign determined by the cobar convention. The right side
is the constant cocycle \eqref{top:beta-cocycle}. Thus the bottom
image of $\beta$ has zero image on $E_2$.

There is no possible higher-filtration homotopy value hidden by
this boundary. The spectrum $DT\wedge Y$ is connective and its
coefficient weights are nonnegative. In stem $pq-2$, its
$q$-sparse terms have
\[
 (s,w)=(2+qj,p+j).
\]
For $j\geq1$ the inequality $w\geq s$ would give
$p-2\geq(q-1)j$, which is impossible. Naturality places the
image of $\beta$ in filtration at least two, and the boundary
just found raises it to filtration at least three. All later
associated-graded groups are zero, so convergence with finite
filtration makes the actual image zero.

The bottom beta image is therefore zero, and finite duality makes
$f$ itself nullhomotopic.
\end{proof}

\subsection{The averaged lift and the page on which it vanishes}

\begin{proposition}[Averaging and annihilation]
\label{top:annihilation}
If $\gamma\in\pi_{qW-5}Y$ represents a class of
$E_2^{5,W}(Y)$, then
\begin{equation}
\label{top:beta-annihilation}
 \beta^K\gamma=0,\qquad K=p(p-2)+1.
\end{equation}
Only the binary two-sided unit identities for $\mu$ are needed.
\end{proposition}

\begin{proof}
Lemma~\ref{top:small-stems} gives $\alpha\gamma=0$ in precisely
the stem $q(W+1)-6$. Proposition~\ref{top:extended-power-relation}
gives $\alpha\beta^p=0$ after applying the unit to $Y$.
By \eqref{top:lift-criterion}, choose lifts
\[
 \widetilde\gamma:S^{qW-5}\longrightarrow J\wedge Y,
 \qquad
 \widetilde b:S^{pd}\longrightarrow J\wedge Y
\]
of $\gamma$ and $\eta\beta^p$, respectively.
Fix any parenthesization of the $n$-fold iteration of $\mu$.
Smash one copy of $\widetilde\gamma$ with $p-2$ copies of
$\widetilde b$, collect the $J$ factors, and apply that chosen
iteration to the $Y$ factors. This gives
\[
 L:S^{qW-5+(p-2)pd}\longrightarrow J^{\wedge n}\wedge Y.
\]
Its top projection is $\beta^{p(p-2)}\gamma$.
To verify this without an associativity assumption, for any
sphere scalar $a$ use the two unit homotopies to identify
\[
 \mu(\eta a,x)=a x,\qquad
 \mu(x,\eta a)=(-1)^{|a||x|}a x.
\]
The scalar $\beta^p$ is even. In the chosen binary tree, every
subtree without $\gamma$ is a scalar times the unit, and the
unique subtree containing $\gamma$ is the product of its scalars
times $\gamma$. This proves the asserted top value by induction
on the tree.

Now project to the symmetric retract after those $Y$ factors
have already been multiplied:
\[
 L_T=(\operatorname{pr}\wedge1_Y)L:
 S^{qW-5+(p-2)pd}\longrightarrow T\wedge Y.
\]
For every permutation $\sigma$ of the $J$ factors, naturality of
the symmetry and the zero-dimensional top sphere give an equality
of stable maps
\[
 \varepsilon^{\wedge n}\sigma=\varepsilon^{\wedge n}.
\]
Consequently
\begin{equation}
\label{top:averaged-top-map}
 (\varepsilon_T\wedge1_Y)L_T
 = (\varepsilon^{\wedge n}e\wedge1_Y)L
 = (\varepsilon^{\wedge n}\wedge1_Y)L
 =\beta^{p(p-2)}\gamma.
\end{equation}
No permutation of the multiplied $Y$ factors occurs in this
identity.

The actual null map $f$ of Proposition~\ref{top:symmetric-null-map}
now gives the following homotopy-commutative diagram:
\begin{equation}
\label{top:null-map-diagram}
\begin{array}{ccc}
\Sigma^d(T\wedge Y)&\xrightarrow{\ f\wedge1_Y\ }&Y\wedge Y\\[2pt]
\mathllap{\scriptstyle\Sigma^d(\varepsilon_T\wedge1_Y)\,}\big\downarrow
&&\big\downarrow\mathrlap{\,\scriptstyle\mu}\\[2pt]
\Sigma^dY&\xrightarrow{\ \beta\wedge1_Y\ }&Y.
\end{array}
\end{equation}
Commutativity is the left unit identity for $\mu$ and the
definition of $f$. Its top arrow is null, so composing the
diagram with $\Sigma^dL_T$ and using
\eqref{top:averaged-top-map} proves
$\beta^{p(p-2)+1}\gamma=0$.
\end{proof}

\begin{proposition}[Disappearance before the last differential]
\label{top:early-disappearance}
Every class in $E_2^{qp+1,u_4+p}(Y)$ is zero on the
Adams--Novikov page $E_{qp+1}$.
\end{proposition}

\begin{proof}
Let $\theta\in E_2^{5,W}(Y)$. By Lemma~\ref{top:small-stems},
it is a permanent cycle represented by a homotopy class $\gamma$.
The sphere class $\beta$ is also permanent. The pagewise product
$A^K\theta$ therefore supports no nonzero outgoing differential;
if either factor has zero image on a page, the product is already
zero. By Proposition~\ref{top:annihilation} its filtered homotopy
product is zero. Product compatibility and finite convergence
then imply that $A^K\theta$ is zero at $E_\infty$.

If it is nonzero at $E_2$, it must consequently be hit by an
incoming differential. Its cohomological filtration is
\[
 S=5+2K=2p^2-4p+7.
\]
An incoming differential of length $r$ has source filtration
$S-r\geq0$, so $r\leq S$. This bounds the page of disappearance,
rather than merely its eventual abutment. The exact comparisons
are
\begin{equation}
\label{top:strict-page-bound}
 qp+1-S=2p-6>0,\qquad N-K=p-3>0.
\end{equation}
Numerically,
\[
 K=1{,}016{,}064,\quad N=1{,}017{,}070,\quad
 S=2{,}032{,}133<qp+1=2{,}034{,}145.
\]
Thus $A^K\theta$ is zero before the page of the prospective
$d_{qp+1}(v_4)$. Multiplication by the further power $A^{N-K}$ is
defined on every page by the sphere permanent cycle $\beta$;
it shows that $A^N\theta$ is zero there as well.
The surjection \eqref{top:An-surjection} accounts for every
$E_2$ class of the target, proving the assertion.
\end{proof}

\subsection{The realizing cofiber}\label{top:main-construction}

All possible targets of a differential from $v_4$ have now been
eliminated on the required pages. We finish by producing the ordinary
finite spectrum and identifying its homology as a comodule.

\begin{proof}[Proof of Theorem~\ref{intro:main}]
The only possible positive differentials from $v_4$ have target
$(s,w)=(qk+1,u_4+k)$. Proposition~\ref{may:first-differential}
rules out $k=1$. Lemma~\ref{may:target-bound} and
Table~\ref{may:remaining-table} leave, for $k\geq2$, only
$k=2,p-1,p$. Corollary~\ref{frob:vanishing} eliminates the
first two target groups already at $E_2$, and
Proposition~\ref{top:early-disappearance} makes the last target
zero before its possible differential. No differential can hit
filtration zero. Hence $v_4\in E_2^{0,u_4}(Y)$ is permanent.

The edge homomorphism of the convergent Adams--Novikov spectral
sequence supplies a map
$a:S^{qu_4}\to Y$ whose $BP$-Hurewicz image is exactly $v_4$.
Extend it by the existing binary multiplication to
\[
 \widetilde v_4:\Sigma^{qu_4}Y
 \xrightarrow{\ a\wedge1_Y\ }Y\wedge Y
 \xrightarrow{\ \mu\ }Y,
\]
and take its cofiber $X$.
The right unit identity shows that the bottom generator is sent
to $v_4$. Since $BP_*Y=R/I_4$ is cyclic as an $R$-module, this
determines the entire induced map as multiplication by $v_4$.
That element is a nonzerodivisor in
$R/I_4=\mathbb F_p[v_4,v_5,\ldots]$. The homology exact sequence
of the cofiber therefore gives the short exact sequence
\[
 0\longrightarrow\Sigma^{qu_4}R/I_4
 \xrightarrow{\ v_4\ }R/I_4
 \longrightarrow BP_*X\longrightarrow0.
\]
The composite bottom map $S^0\to Y\to X$ induces the quotient
map $R\to R/(I_4,v_4)$. It is a comodule map and is surjective,
so it identifies the induced coaction with the canonical quotient
coaction, and identifies its generator in degree zero.

Finally, $X$ is the cofiber of a map between finite $p$-local
spectra. It is finite $p$-local in the sense defined in the introduction;
equivalently, the finite cell construction can be realized by a
finite spectrum before localization by clearing prime-to-$p$
denominators in its attaching maps, as in
Lemma~\ref{low:finite-localization}. More explicitly, the mapping cone
adds to the cells of $Y$ their translates by $qu_4+1$.
Together with Equation~\eqref{low:cell-dimensions}, this gives a finite
$p$-local CW model for $X$ with one cell in each dimension
\[
 \sum_{i\in S}(qu_i+1)=\sum_{i\in S}(2p^i-1),
 \qquad S\subseteq\{0,1,2,3,4\}.
\]
In particular, this model has $32$ cells.
\end{proof}
